% AdvancedThermalQuality20261023-v1
% Editable mathematical diagrams: index structure, orthogonal directions and boundaries.
\providecommand{\ThermalFigure}[3]{\begin{qualityfigure}\begin{tikzpicture}[>=Latex,font=\small,x=.8cm,y=.8cm]#1\end{tikzpicture}\figcap{#2}\label{#3}\end{qualityfigure}}
\newcommand{\AQEpsilon}{\ThermalFigure{
 \draw[thick](0,0)grid[step=.7](2.1,2.1);\foreach\x/\y in {.35/1.75,1.05/1.05,1.75/.35}{\node at(\x,\y){$1$};}
 \node[left]at(0,1.75){$[1]_r$};\node[left]at(0,1.05){$[2]_r$};\node[left]at(0,.35){$[3]_r$};
 \node[right,align=left]at(3,1.2){从单位矩阵依次选行 $i,j,k$\\未标格的元素为零};
}{每行是一个标准基矢的分量，指标重复与行重复相对应。}{advA3:quality:epsilon-q}}
\newcommand{\AAEpsilon}{\ThermalFigure{
 \node[draw,align=center](A)at(1,1){行序 $1,2,3$\\正定向};\node[draw,align=center](B)at(6,1){行序 $2,1,3$\\交换两行};
 \draw[<->,thick](A)--node[above]{一次换位}(B);\node[below]at(3.5,-.5){行列式 $+1\longleftrightarrow-1$；两行重复时为零};
}{定向改变对应全反对称符号；缩并后的 Kronecker 符号表达相同的指标配对。}{advA3:quality:epsilon-a}
图~\ref{advA3:quality:epsilon-a}中一次换位翻转定向，循环置换包含两次换位而保号。缩并 $\varepsilon_{ijk}\varepsilon_{imn}$ 时，若 $j=k$ 两边均为零；若 $j\ne k$，第一个因子只允许互补的一个 $i$，第二个因子再检验 $(m,n)$ 是同序还是反序，从而得到两个 $\delta$ 乘积之差。
}
\newcommand{\AQJacobi}{\ThermalFigure{
 \foreach\x/\y/\lab in {3/0/u,1/2.5/v,-1.5/1.2/w}{\draw[->,thick](0,0)--(\x,\y)node[above]{$\vect\lab$};}\fill(0,0)circle(2pt)node[below]{$O$};
 \node[right]at(4,1.5){同起点矢量的投影示意，不预设三矢量共面};
}{三个矢量在双重叉积中循环交换，括号的次序须保留。}{advA3:quality:jacobi-q}}
\newcommand{\AAJacobi}{\ThermalFigure{
 \node[draw,align=left]at(4,2){沿 $\vect u$：$-(\vect v\cdot\vect w)\vect u+(\vect w\cdot\vect v)\vect u$};
 \node[draw,align=left]at(4,.7){沿 $\vect v$：$(\vect u\cdot\vect w)\vect v-(\vect w\cdot\vect u)\vect v$};
 \node[draw,align=left]at(4,-.6){沿 $\vect w$：$-(\vect u\cdot\vect v)\vect w+(\vect v\cdot\vect u)\vect w$};
}{展开后三个方向的系数分别成对抵消，不需要三个矢量线性无关。}{advA3:quality:jacobi-a}
图~\ref{advA3:quality:jacobi-a}把六项按同一输出矢量分组。每组的两个标量内积相等、符号相反，故逐组为零；证明不依赖三矢量能否组成一组基。这也说明应使用双重叉积恒等式，不能将叉积误当作满足结合律的乘法。
}
\newcommand{\AQDerivations}{\ThermalFigure{
 \node[draw](W)at(0,1){$\vect w$};\node[draw](V)at(4,2){$D_{\vect v}\vect w$};\node[draw](U)at(4,0){$D_{\vect u}\vect w$};
 \draw[->](W)--node[above]{$D_{\vect v}$}(V);\draw[->](W)--node[below]{$D_{\vect u}$}(U);
 \node[right]at(7,1){继续复合并比较两种次序};
}{算子复合中右侧算子先作用；两条路径的差是对易子。}{advA3:quality:derivation-q}}
\newcommand{\AADerivations}{\ThermalFigure{
 \node[draw,align=center](A)at(1,1){$D_{\vect u}D_{\vect v}\vect w$\\$-D_{\vect v}D_{\vect u}\vect w$};
 \node[draw](B)at(8,1){$(\vect u\times\vect v)\times\vect w$};\draw[->,thick](A)--node[above]{雅可比重排}(B);
 \node[below]at(4.5,-.4){同一等式对每个输入 $\vect w$ 成立，才得到算子恒等式};
}{把对易子先落实到任意测试矢量，避免将形式符号当作交换乘法。}{advA3:quality:derivation-a}
图~\ref{advA3:quality:derivation-a}中的右端是一个新的叉积算子对 $\vect w$ 的作用。由于 $\vect w$ 任意，才可去掉测试输入得到算子等式。导子性质则将 $D_{\vect u}$ 分别作用于楔积的两个因子，输出仍为同一矢量空间中的矢量。
}
\newcommand{\AQSensors}{\ThermalFigure{
 \draw[thick](0,0)grid[step=.8](2.4,2.4);\fill[DeepBlue!12](0,0)--(2.4,2.4)--(2.4,0)--cycle;
 \node[right,align=left]at(3.2,1.6){二阶例：$T_{ij}=T_{ji}$\\一般 $n$ 阶：所有指标置换均对称};
}{完全对称性把指标顺序不同、出现次数相同的分量识别为同一个自由参数。}{advA3:quality:symmetric-q}}
\newcommand{\AASensors}{\ThermalFigure{
 \foreach\x in {0,.7,2.2,3.7,4.4}{\fill[DeepBlue](\x,0)circle(3pt);}\draw[very thick](1.4,-.5)--(1.4,.5) (2.9,-.5)--(2.9,.5);
 \node[below]at(.35,-.7){$r_1=2$};\node[below]at(2.2,-.7){$r_2=1$};\node[below]at(4.05,-.7){$r_3=2$};
 \node[above]at(2.2,1.1){五颗星、两块隔板：五阶三维的一种分量类型};
}{允许相邻隔板，表示某类指标可以出现零次；总位置数为 $n+m-1$。}{advA3:quality:symmetric-a}
图~\ref{advA3:quality:symmetric-a}中的点表示指标出现次数而不是不同粒子。完全对称性使每一组计数 $(r_1,\ldots,r_m)$ 只对应一个独立分量，因此从 $n+m-1$ 个位置选 $m-1$ 个放隔板即可。若张量只对部分指标对称，这一计数不再适用。
}
\newcommand{\AQExterior}{\ThermalFigure{
 \foreach\x/\y/\lab in {3/0/1,0/2.5/2,-1.8/-1.2/3}{\draw[->,thick](0,0)--(\x,\y)node[right]{$\vect e_{\lab}$};}
 \node[right]at(4,1.5){各阶基由不同基矢的楔积构成};
}{外代数同时含零阶标量、向量及更高阶有向面积与体积元。}{advA3:quality:exterior-q}}
\newcommand{\AAExterior}{\ThermalFigure{
 \foreach\x/\w/\lab in {0/1/{1},1/3/{e_1,e_2,e_3},4/3/{e_1\wedge e_2,e_1\wedge e_3,e_2\wedge e_3},7/1/{e_1\wedge e_2\wedge e_3}}{\draw[thick](\x,0)rectangle(\x+\w,1);}
 \node at(.5,.5){$1$};\node at(2.5,.5){$\Lambda^1$：3};\node at(5.5,.5){$\Lambda^2$：3};\node at(7.5,.5){1};
 \node[below]at(4,-.6){三维：$1+3+3+1$；每个基矢可选或不选};
}{每个递增指标子集唯一对应一个楔积基元，重复指标则为零。}{advA3:quality:exterior-a}
图~\ref{advA3:quality:exterior-a}按阶数分组计数，最左的一份必须包括空子集对应的标量单位。每个基矢独立有“选、不选”两种选择，给出 $2^n$ 个子集；固定选 $l$ 个时才得到 $\binom nl$，两种计数相互核验。
}
\newcommand{\AQHodge}{\ThermalFigure{
 \fill[DeepBlue!10](-2,-.5)--(1,-.5)--(2,.5)--(-1,.5)--cycle;\draw[thick](-2,-.5)--(1,-.5)--(2,.5)--(-1,.5)--cycle;
 \draw[->,thick](0,0)--(0,2.6)node[right]{$\vect v$};\draw[->,DeepBlue](0,0)--(2,2)node[right]{$\vect u$};\node[below]at(0,-1){$\star\vect v$ 所在有向平面};
}{三维欧氏空间中的矢量与有向面积对偶，采用固定正定向。}{advA3:quality:hodge-q}}
\newcommand{\AAHodge}{\ThermalFigure{
 \draw[thick](-2,0)--(3,0);\node[below]at(-1,-.2){$\star\vect v$ 的平面};\draw[->](0,0)--(0,2.5)node[left]{$\vect v$};
 \draw[->,DeepBlue,thick](0,0)--(2,2)node[right]{$\vect u$};\draw[dashed](2,2)--(0,2);\draw[->,thick](0,0)--(0,2)node[left]{$u_\parallel$};
}{楔积产生有向体积，只保留 $\vect u$ 沿平面法向的分量。}{advA3:quality:hodge-a}
图~\ref{advA3:quality:hodge-a}中切向分量与面积元线性相关，其楔积为零，法向分量则给出 $|u_\parallel||v|$ 的有向体积。最后的霍奇对偶将三阶体积元变成标量，结果正是内积。若改变空间定向，两次对偶的符号须一起改变，不能只更改一次。
}
\newcommand{\AQRadialHessian}{\ThermalFigure{
 \foreach\r in {.7,1.3,2}{\draw[TextGray](0,0)circle(\r);}\fill(0,0)circle(2pt)node[below]{$O$};\fill[DeepBlue](1.414,1.414)circle(2pt)node[right]{$P$};\draw[->,thick](0,0)--(1.414,1.414)node[midway,above left]{$\vect r$};
 \node[right]at(3.5,1){$f(r)$ 的等值面为同心球面};
}{径向函数只通过 $r$ 依赖各坐标，导数仍须按坐标分量逐一求取。}{advA3:quality:radial-hessian-q}}
\newcommand{\AARadialHessian}{\ThermalFigure{
 \draw[thick](0,0)circle(2);\fill(2,0)circle(2pt)node[below right]{$P$};\draw[->,DeepBlue,thick](2,0)--(3.3,0)node[right]{径向};\draw[->,thick](2,0)--(2,1.4)node[right]{切向};\draw[dashed](0,0)--(2,0);
 \node[right,align=left]at(4.8,1.5){径向本征值 $n(n-1)r^{n-2}$\\两个切向本征值 $nr^{n-2}$};
}{Hessian 在径向与两切向分量上分别作用，其迹由三个本征值相加。}{advA3:quality:radial-hessian-a}
图~\ref{advA3:quality:radial-hessian-a}将局部增量分为沿 $\vect r$ 与垂直 $\vect r$ 的方向。由 $\partial_i r=x_i/r$ 再求一次导数可得 $\partial_i\partial_jr^n=nr^{n-2}\delta_{ij}+n(n-2)r^{n-4}x_ix_j$；切向增量与 $\vect r$ 点乘为零，直接读出两个本征值。所有这些点态表达式都排除原点的分布项。
}
\newcommand{\AQGreen}{\ThermalFigure{
 \draw[thick,fill=DeepBlue!7](0,0)ellipse(2.5 and 1.5);\node at(0,0){$V:\ \phi,\psi$};\fill(2.5,0)circle(2pt);\draw[->,thick](2.5,0)--(3.7,0)node[right]{$\vect n$};\node[above]at(0,2){分片光滑边界 $\partial V$};
}{边界法向取向外，曲面元为 $\dd\vect S=\vect n\dd S$。}{advA3:quality:green-q}}
\newcommand{\AAGreen}{\ThermalFigure{
 \draw[fill=DeepBlue!7,thick](-3,-1.3)rectangle(3,1.3);\draw[thick](0,-1.3)--(0,1.3);\node at(-1.5,.6){$V_1$};\node at(1.5,.6){$V_2$};
 \draw[->,DeepBlue,thick](0,0)--(1,0)node[below]{$\vect n_1$};\draw[->,thick](0,0)--(-1,0)node[below]{$\vect n_2$};
 \draw[->](3,.7)--(4.2,.7)node[right]{外边界法向};
}{相邻子体积的内部通量相消，只有整个区域的外边界保留下来。}{advA3:quality:green-a}
图~\ref{advA3:quality:green-a}中内部面的两法向相反，同一光滑通量 $\phi\nabla\psi$ 的贡献因而抵消。第一恒等式来自这一通量的散度；交换 $\phi,\psi$ 后相减，两个相同的梯度内积也抵消，从而得到第二恒等式。场若含奇点，应先挖去奇点小邻域，并另外保留其内边界通量。
}

\section{矢量分析}

\subsection{指标语言与三维矢量}

指标记号把成组分量、求和和缩并统一到一套语法中。例如
\[
 a_i=(1,2,3),\qquad b_i=(4,5,6),\qquad
 c_{ij}=a_ib_j,\qquad a_ib_i=32.
\]
只出现一次的 $i,j,k$ 是\emph{自由指标}，等式两边必须逐一对应；同一项中恰好出现两次的指标是\emph{哑指标}，按爱因斯坦求和约定对其全部取值求和。不同自由指标独立取值，三维中含两个自由指标的量有 $3^2$ 个分量。哑指标可改名，却不能与同项已有自由指标冲突。

从四个最小表达式可以引入这套语言：
\[
 a_i=(1,2,3),\qquad b_i=(4,5,6),
\]
\[
 c_{ij}=a_ib_j=
 \begin{pmatrix}4&5&6\\8&10&12\\12&15&18\end{pmatrix},
 \qquad d=a_ib_i=1\cdot4+2\cdot5+3\cdot6=32.
\]
前两式各含一个自由指标，实际各代表三个分量；第三式的 $i,j$ 独立变化，代表九个分量；第四式的 $i$ 在同一单项式中出现两次，因而是哑指标，只代表一次求和。由此得到三条必须始终遵守的语法规则：

\begin{enumerate}
 \item 同一单项式中，一个指标只能出现一次（自由）或两次（缩并），不能出现三次，例如 $a_ib_ic_i$ 没有定义；
 \item 等式两边的自由指标必须相同，例如 $A_{ij}=B_{ik}$ 不成立，因为右侧还含未缩并的 $k$；
 \item 不同单项式中的哑指标可以独立改名，例如
 $a_ib_i+c_jd_j=a_kb_k+c_ld_l$，但不能把它改成与同项自由指标冲突的名字。
\end{enumerate}

这些规则并非排版习惯，而是保证一个短式能够无歧义地还原成分量方程。比如
$A_{ij}x_j=b_i$ 代表三条方程，每条方程内部对 $j$ 求和；而
$A_{ij}x_i=b_j$ 则是另一个转置后的线性映射。指标的上下位置在一般坐标系中还区分逆变与协变分量；本节限于欧氏正交归一基，可用度量 $\delta_{ij}$ 升降指标，所以数值上 $x_i=x^i$，但概念上仍保留位置区别。

\begin{advdefinition}{克罗内克符号}\label{advA3:def:kronecker}
\[
 \delta_{ij}=\begin{cases}1,&i=j,\\0,&i\ne j.
 \end{cases}
\]
它是欧氏正交归一基中的单位张量分量。
\end{advdefinition}

\begin{advdefinition}{列维--奇维塔符号}\label{advA3:def:levi-civita}
三维中
\[
 \varepsilon_{ijk}=\begin{cases}
 +1,&(ijk)\text{ 是 }(123)\text{ 的偶置换},\\
 -1,&(ijk)\text{ 是 }(123)\text{ 的奇置换},\\
 0,&i,j,k\text{ 中有重复}.
 \end{cases}
\]
\end{advdefinition}
$\delta_{ij}$ 缩并时起换指标作用，例如
$a_ib_j\delta_{ij}=a_ib_i$、
$\delta_{ij}\delta_{kl}A_{ikm}=A_{jlm}$。

把换指标过程逐项写开可看得更清楚。对固定自由指标 $i$，
\[
 \delta_{ij}v_j=\delta_{i1}v_1+\delta_{i2}v_2+\delta_{i3}v_3=v_i.
\]
因此 $\delta_{ij}$ 是单位映射的分量。它的迹为
$\delta_{ii}=3$；而 $\delta_{ij}\delta_{jk}=\delta_{ik}$ 表明单位映射复合后仍是自身。列维--奇维塔符号则记录三维空间的取向：只有三个指标互异时才非零，循环置换 $(123)\to(231)\to(312)$ 保持正号，任意一次对换改变符号。

\begin{advtheorem}{列维--奇维塔缩并恒等式}\label{advA3:thm:epsilon-identities}
\[
 \varepsilon_{ijk}\varepsilon_{ijk}=6,
 \qquad
 \varepsilon_{ijk}\varepsilon_{ijn}=2\delta_{kn},
\]
\[
 \boxed{\varepsilon_{ijk}\varepsilon_{imn}
 =\delta_{jm}\delta_{kn}-\delta_{jn}\delta_{km}},
\]
\[
 \varepsilon_{ijk}\varepsilon_{lmn}
 =\det\!\begin{pmatrix}
 \delta_{il}&\delta_{im}&\delta_{in}\\
 \delta_{jl}&\delta_{jm}&\delta_{jn}\\
 \delta_{kl}&\delta_{km}&\delta_{kn}
 \end{pmatrix}.
\]
\end{advtheorem}
最后一式是两个有向体积形式的内积展开；令 $l=i$ 并缩并，依次得到前三式。

还可以不用预设结论来检查最常用的二指标缩并。固定 $j,k,m,n$，对 $i$ 求和。若 $j=k$，左边因第一个 $\varepsilon$ 为零，右边也成为
$\delta_{jm}\delta_{jn}-\delta_{jn}\delta_{jm}=0$。若 $j\ne k$，只有唯一一个与 $j,k$ 互补的 $i$ 可能贡献；第二个 $\varepsilon_{imn}$ 也只有当 $(m,n)$ 等于 $(j,k)$ 或 $(k,j)$ 时非零，正好分别对应右边的 $+1$ 与 $-1$。这给出独立的分量核验。

\begin{advexamplebox}{用行列式构造列维--奇维塔符号}\label{advA3:ex:epsilon-determinant}
令 $[i]_r$、$[i]_c$ 分别为三阶单位矩阵的第 $i$ 行和第 $i$ 列。
\begin{enumerate}
 \item 证明
 $\varepsilon_{ijk}=\det([i]_r,[j]_r,[k]_r)^{\mathsf T}$；
 \item 由此证明 $\varepsilon_{ijk}$ 全反对称；
 \item 证明定理\ref{advA3:thm:epsilon-identities}。
\end{enumerate}

\AQEpsilon
\end{advexamplebox}
\begin{advsolution}
\AAEpsilon

若 $i,j,k$ 有重复，行列式有两行相同而为零；若互异，这三行恰是单位阵三行的一个置换，行列式等于置换符号，正是 $\varepsilon_{ijk}$。交换任意两个指标等于交换两行，故符号反号。

对两个行列式用柯西--比内公式，
\[
 \varepsilon_{ijk}\varepsilon_{lmn}
 =\det([i]_r,[j]_r,[k]_r)^{\mathsf T}
  \det([l]_c,[m]_c,[n]_c),
\]
矩阵乘积的九个元素就是相应的 $\delta$，得到定理最后一式。再令
$l=i$ 求和；只有与固定 $(j,k)$ 互补的一个 $i$ 有贡献，展开行列式即得
$\delta_{jm}\delta_{kn}-\delta_{jn}\delta_{km}$。继续令 $m=j$、$n=k$ 缩并，分别得到 $2\delta_{kn}$ 和6。
\end{advsolution}

\begin{advdefinition}{三维欧氏矢量空间}\label{advA3:def:vector-space}
取线性无关基 $\vect e_1,\vect e_2,\vect e_3$，三维实矢量空间为
\[
 \mathbb R^3=\{x^i\vect e_i\mid x^i\in\mathbb R\}.
\]
元素 $\vect x=x^i\vect e_i$ 称为矢量；$x^i$ 只是该矢量在选定基下的分量。抽象指标写法 $x^a$ 强调“这是一个矢量”而不指定基。
\end{advdefinition}

这里要严格区分三个对象：$\vect e_i$ 是基矢，$x^i$ 是三个实数分量，$x^a$ 是与选基无关的抽象矢量。若换基
$\vect e'_i=\vect e_jM^j{}_i$，其中矩阵 $M$ 的第 $i$ 列为新基矢 $\vect e'_i$ 在旧基中的分量。同一个几何矢量的坐标列必须满足 $x=Mx'$，因而反向变化，
\[
 x'^i=(M^{-1})^i{}_jx^j,
 \qquad x'^i\vect e'_i=x^j\vect e_j.
\]
所以“把三个数排成一列”只有连同它们所依赖的基才表示矢量。基在定义线性空间时可以先视为“忘却结构的符号”：暂不赋予长度、角度或方向，随后再由点积与取向结构赋予欧氏几何含义。

抽象指标 $a,b,c$ 不取 $1,2,3$ 中的数值，它只标明对象所在的张量空间。例如 $u^av^b$ 是二阶张量，$u^av_a$ 是缩并得到的标量；把抽象指标换成具体指标并补上基，才得到可计算的分量式。后续若不讨论换基，通常仍用粗体矢量记号提高可读性，但每一个矢量恒等式都可以退回指标式审查。

\begin{advdefinition}{点积}\label{advA3:def:dot}
点积是双线性、对称、正定映射 $\mathbb R^3\times\mathbb R^3\to\mathbb R$，在正交归一基中
$\vect e_i\cdot\vect e_j=\delta_{ij}$，所以
$\vect x\cdot\vect y=x^iy^j\delta_{ij}=x^iy^i$。
\end{advdefinition}

\begin{advdefinition}{叉积}\label{advA3:def:cross}
在定向三维欧氏空间中，叉积是双线性映射
$\mathbb R^3\times\mathbb R^3\to\mathbb R^3$，由
\[
 \vect e_i\times\vect e_j=\varepsilon_{ijk}\vect e_k
\]
确定方向与大小。
\end{advdefinition}

\begin{advtheorem}{点积交换、叉积反交换}\label{advA3:thm:dot-cross-symmetry}
\[
 \vect u\cdot\vect v=\vect v\cdot\vect u,
 \qquad
 \vect u\times\vect v=-\vect v\times\vect u.
\]
\end{advtheorem}
证明只需分别使用 $\delta_{ij}=\delta_{ji}$ 与
$\varepsilon_{ijk}=-\varepsilon_{jik}$。

分量计算同时还原了熟悉的坐标公式：
\[
 \vect x\cdot\vect y=x_1y_1+x_2y_2+x_3y_3,
\]
\[
 \vect x\times\vect y=
 (x_2y_3-x_3y_2)\vect e_1
 +(x_3y_1-x_1y_3)\vect e_2
 +(x_1y_2-x_2y_1)\vect e_3.
\]
由 Cauchy--Schwarz 不等式可定义夹角
$\cos\theta=(\vect x\cdot\vect y)/(|\vect x||\vect y|)$；而叉积大小满足
\[
 |\vect x\times\vect y|^2
 =x_iy_jx_my_n\varepsilon_{ijk}\varepsilon_{mnk}
 =|\vect x|^2|\vect y|^2-(\vect x\cdot\vect y)^2,
\]
所以 $|\vect x\times\vect y|=|\vect x||\vect y|\sin\theta$，方向由空间取向和右手规则确定。

\begin{advdefinition}{三重标积}\label{advA3:def:triple-scalar}
\[
 \mathcal V(\vect u,\vect v,\vect w)
 =(\vect u\times\vect v)\cdot\vect w
 =u_iv_jw_k\varepsilon_{ijk}.
\]
它是三矢量张成平行六面体的有向体积。
\end{advdefinition}

\begin{advtheorem}{三重标积的刻画}\label{advA3:thm:triple-properties}
$\mathcal V$ 对三个变量分别线性，任意两个输入相等时为零，循环置换不变，并满足
$\mathcal V(\vect e_1,\vect e_2,\vect e_3)=1$。反过来，这四条性质唯一确定三重标积。
\end{advtheorem}
展开到基矢后，交错性把所有含重复指标的项消去，正规性和循环性把其余六项化成置换符号，因此必为
$u_iv_jw_k\varepsilon_{ijk}$。

用分量矩阵写就是
\[
 \mathcal V(\vect u,\vect v,\vect w)=
 \det\!\begin{pmatrix}
 u_1&u_2&u_3\\v_1&v_2&v_3\\w_1&w_2&w_3
 \end{pmatrix}.
\]
交换两行改变取向，所以反号；若三个矢量共面，则三行线性相关，体积为零。反过来，只要这一行列式非零，三个矢量就构成一组基。三重标积因而同时承担几何体积、线性无关判据和坐标变换 Jacobian 三种角色。

\begin{advtheorem}{三重矢积公式}\label{advA3:thm:triple-vector}
\[
 (\vect u\times\vect v)\times\vect w
 =(\vect u\cdot\vect w)\vect v-(\vect v\cdot\vect w)\vect u,
\]
\[
 \vect u\times(\vect v\times\vect w)
 =(\vect u\cdot\vect w)\vect v-(\vect u\cdot\vect v)\vect w.
\]
\end{advtheorem}
例如第一式的第 $l$ 分量为
$u_iv_jw_k\varepsilon_{ijm}\varepsilon_{mkl}$；用定理\ref{advA3:thm:epsilon-identities}化为
$u_iw_iv_l-v_iw_iu_l$，即得。

\begin{advexamplebox}{叉积的雅可比恒等式}\label{advA3:ex:jacobi-cross}
证明
\[
 (\vect u\times\vect v)\times\vect w
 +(\vect w\times\vect u)\times\vect v
 +(\vect v\times\vect w)\times\vect u=0.
\]

\AQJacobi
\end{advexamplebox}
\begin{advsolution}
\AAJacobi

逐项使用定理\ref{advA3:thm:triple-vector}：
\[
 (u\cdot w)v-(v\cdot w)u
 +(w\cdot v)u-(u\cdot v)w
 +(v\cdot u)w-(w\cdot u)v=0.
\]
每一方向的两项因点积交换律相消。
\end{advsolution}

\begin{advexamplebox}{叉积导子与对易子}\label{advA3:ex:cross-derivation}
定义 $D_{\vect u}(\vect v)=\vect u\times\vect v$，并令
$[L_1,L_2]=L_1L_2-L_2L_1$。证明
\[
 [D_{\vect u},D_{\vect v}]=D_{\vect u\times\vect v},
 \qquad
 D_{\vect u}(\vect v\times\vect w)
 =(D_{\vect u}\vect v)\times\vect w
 +\vect v\times(D_{\vect u}\vect w).
\]

\AQDerivations
\end{advexamplebox}
\begin{advsolution}
\AADerivations

对任意 $\vect w$，雅可比恒等式重排为
\[
 \vect u\times(\vect v\times\vect w)
 -\vect v\times(\vect u\times\vect w)
 =(\vect u\times\vect v)\times\vect w,
\]
即第一式。第二式直接展开右边并用三重矢积：
\[
 (u\times v)\times w+v\times(u\times w)
 =(u\cdot w)v-(v\cdot w)u+(v\cdot w)u-(u\cdot v)w,
\]
等于 $u\times(v\times w)$。
\end{advsolution}

\subsection{张量与外代数}

\begin{advdefinition}{张量与张量积}\label{advA3:def:tensor-product}
二阶张量空间是
\[
 \mathbb R^3\otimes\mathbb R^3
 =\{T^{ij}\vect e_i\vect e_j\},
\]
张量积由双线性和
$\vect e_i\otimes\vect e_j=\vect e_i\vect e_j$ 定义。高阶张量用更多有序指标表示；零阶和一阶张量分别是标量和矢量。
\end{advdefinition}
张量积一般不交换：$u^av^b$ 与 $v^au^b$ 的分量顺序不同。各阶张量组成分级代数
$T(\mathbb R^3)=\bigoplus_{n\ge0}T^n(\mathbb R^3)$。

二阶张量在给定基下有九个分量，
\[
 \mathsf T=T^{ij}\vect e_i\otimes\vect e_j.
\]
它可以看成双线性型，也可以在用度量识别矢量和对偶矢量后看成线性映射。二者分量写法相近，但指标作用不同，因而不应仅凭“像矩阵”就随意交换次序。三个矢量的张量积
$u^av^bw^c$ 属于三阶空间；一般有
\[
 T(\mathbb R^3)=T^0\oplus T^1\oplus T^2\oplus\cdots,
 \qquad T^l\otimes T^m\subset T^{l+m}.
\]
标量是零阶张量，矢量是一阶张量。这个分级解释了为什么张量积会增加指标，而缩并会减少两个指标。

\begin{advdefinition}{张量转置}\label{advA3:def:tensor-transpose}
$t_{pq}$ 交换张量第 $p,q$ 个指标。例如
$(t_{12}T)^{ab}=T^{ba}$；矩阵转置只是二阶张量的这一特例。
\end{advdefinition}

\begin{advdefinition}{张量缩并}\label{advA3:def:contraction}
缩并 $c_{pq}$ 把第 $p,q$ 个一上一下指标相等并求和，从而使阶数降低2。例如
对未缩并的二阶张量 $T^a{}_b=u^av_b$，有 $c_{12}T=T^a{}_a=u^av_a=\vect u\cdot\vect v$。二阶混合张量的缩并就是迹，不能把原来两个不同矢量误写成同一矢量的模平方。
\end{advdefinition}
抽象指标能无歧义地区分矩阵乘法的不同转置组合：
$A^{ab}B_{bc}$、$A^{ab}B_{cb}$、$A^{ba}B_{bc}$ 分别代表不同缩并。双点积也是两个指标同时缩并。

更具体地，把 $A,B$ 的分量看作矩阵，四种常见乘积分别为
\[
 A^{ab}B_{bc}\leftrightarrow AB,\qquad
 A^{ab}B_{cb}\leftrightarrow AB^{\mathsf T},
\]
\[
 A^{ba}B_{bc}\leftrightarrow A^{\mathsf T}B,\qquad
 A^{ba}B_{cb}\leftrightarrow A^{\mathsf T}B^{\mathsf T}.
\]
而
$A^{ab}B_{ba}=\operatorname{tr}(AB)$ 与
$A^{ab}B_{ab}=\operatorname{tr}(A^{\mathsf T}B)$
是两种不同的双缩并。抽象指标的价值正在于：每一次转置和缩并都直接写在指标位置上，不依赖另造符号。

\begin{advdefinition}{对称化与反对称化}\label{advA3:def:symmetrization}
\[
 T^{(a_1\cdots a_n)}=\frac1{n!}\sum_{\sigma\in S_n}
 T^{a_{\sigma(1)}\cdots a_{\sigma(n)}},
\]
\[
 T^{[a_1\cdots a_n]}=\frac1{n!}\sum_{\sigma\in S_n}
 \operatorname{sgn}(\sigma)T^{a_{\sigma(1)}\cdots a_{\sigma(n)}}.
\]
两者都是投影算子：$\mathcal S^2=\mathcal S$、$\mathcal A^2=\mathcal A$。
\end{advdefinition}

二阶情形尤其直观：
\[
 T^{(ab)}=\frac12(T^{ab}+T^{ba}),\qquad
 T^{[ab]}=\frac12(T^{ab}-T^{ba}),
\]
\[
 T^{ab}=T^{(ab)}+T^{[ab]}.
\]
对称部分在转置下不变，反对称部分在转置下反号。除以 $n!$ 是必要的：若 $T$ 已完全对称，再做一次对称化必须仍等于 $T$，而不是变成 $n!T$。完全反对称张量只保留指标互异的分量，其余含重复指标的分量全部为零。

\begin{advexamplebox}{对称张量空间的维数}\label{advA3:ex:symmetric-dimension}
证明
\[
 \dim S^n(\mathbb R^3)=\frac{(n+1)(n+2)}2,
\]
并推广到 $S^n(\mathbb R^m)$。

\AQSensors
\end{advexamplebox}
\begin{advsolution}
\AASensors

完全对称分量只由指标1、2、3各出现多少次决定。令出现次数为
$r_1,r_2,r_3\ge0$ 且 $r_1+r_2+r_3=n$，不同非负整数解的个数由插板法为
\[
 \binom{n+2}{2}=\frac{(n+1)(n+2)}2.
\]
$m$ 维时同理数 $r_1+\cdots+r_m=n$ 的非负整数解：
\[
 \boxed{\dim S^n(\mathbb R^m)=\binom{n+m-1}{m-1}}.
\]
\end{advsolution}

三维中 $\dim\Lambda^0=1,\dim\Lambda^1=3,\dim\Lambda^2=3,\dim\Lambda^3=1$。二阶交错基可取
$\vect e_1\wedge\vect e_2$、$\vect e_2\wedge\vect e_3$、$\vect e_3\wedge\vect e_1$；三阶只有体积元
$\vect e_1\wedge\vect e_2\wedge\vect e_3$。

按照本节楔积的归一化，
\[
 \vect e_i\wedge\vect e_j
 =\vect e_i\otimes\vect e_j-\vect e_j\otimes\vect e_i,
\]
而三阶体积元是六个排列项按置换符号求和：
\[
 \begin{aligned}
 \vect e_1\wedge\vect e_2\wedge\vect e_3
 ={}&\vect e_1\vect e_2\vect e_3
 +\vect e_2\vect e_3\vect e_1
 +\vect e_3\vect e_1\vect e_2\\
 &-\vect e_3\vect e_2\vect e_1
 -\vect e_1\vect e_3\vect e_2
 -\vect e_2\vect e_1\vect e_3.
 \end{aligned}
\]
由这些显式展开即可自然过渡到一般交错代数。

\begin{advtheorem}{三维四阶以上交错张量为零}\label{advA3:thm:alternating-zero}
\[
 n\ge4\quad\Longrightarrow\quad \Lambda^n(\mathbb R^3)=\{0\}.
\]
\end{advtheorem}
证明：任意基张量含 $n\ge4$ 个取值只可能为1、2、3的指标，抽屉原理保证至少两个相同。交换这两个位置既不改变基张量，又因全反对称性使其反号，故该项为零。

\begin{advdefinition}{楔积}\label{advA3:def:wedge}
若 $\alpha\in\Lambda^l$、$\beta\in\Lambda^m$，定义
\[
 \alpha\wedge\beta=\frac{(l+m)!}{l!m!}\,
 \mathcal A(\alpha\otimes\beta).
\]
它满足结合律及分次交换律
$\alpha\wedge\beta=(-1)^{lm}\beta\wedge\alpha$。
\end{advdefinition}

分次交换律中的 $(-1)^{lm}$ 来自把 $\alpha$ 的 $l$ 个指标逐个越过 $\beta$ 的 $m$ 个指标，一共需要 $lm$ 次对换。因而两个1-形式反交换，一个1-形式与一个2-形式交换，两个2-形式也交换。特别地，对任意1-形式 $u$，
$u\wedge u=0$。结合律保证可以无括号地写
$u\wedge v\wedge w$，它等于三者张量积的完全反对称部分乘以 $3!$。

\begin{advdefinition}{三维霍奇对偶}\label{advA3:def:hodge}
霍奇算子 $\star:\Lambda^l\to\Lambda^{3-l}$ 线性，并由
\[
 \star1=\vect e_1\wedge\vect e_2\wedge\vect e_3,
 \quad
 \star\vect e_1=\vect e_2\wedge\vect e_3
\]
及循环关系确定；反向又有
$\star(\vect e_2\wedge\vect e_3)=\vect e_1$。三维欧氏正定情形满足 $\star\star=1$。
\end{advdefinition}
由定义
\[
 \star(\vect x\wedge\vect y)=\vect x\times\vect y,
 \qquad
 \star(\vect x\wedge\vect y\wedge\vect z)
 =\mathcal V(\vect x,\vect y,\vect z).
\]
这说明叉积是二矢量交错张量积的霍奇对偶，三重标积是三矢量楔积的霍奇对偶。

逐个基矢核验第一式：
\[
 \star(\vect e_i\wedge\vect e_j)=\varepsilon_{ijk}\vect e_k,
\]
恰与叉积定义相同；由双线性即可推广到任意 $\vect x,\vect y$。再楔入 $\vect z$ 并取对偶，
\[
 \star(\vect x\wedge\vect y\wedge\vect z)
 =x_iy_jz_k\varepsilon_{ijk},
\]
就是三重标积。于是三维叉积并不是一般维数都自然存在的基本乘法，而是“三维二阶交错张量与一阶矢量恰好同维”带来的特殊对偶。

若记体积元
$\Omega=\vect e_1\wedge\vect e_2\wedge\vect e_3$，则内插乘积满足
\[
 \iota_{\vect u}\Omega=\star\vect u,
 \qquad \iota_{\vect v}\iota_{\vect u}\Omega
 =\star(\vect u\wedge\vect v).
\]
也可用记号 $\vect e_i\cdot\Omega=\star\vect e_i$ 表达同一关系，并由此把三重矢积重写为
\[
 (\vect x\times\vect y)\times\vect z
 =\iota_{\vect z}(\vect x\wedge\vect y)
 =(\vect x\cdot\vect z)\vect y-(\vect y\cdot\vect z)\vect x.
\]
这里右式也再次说明三重矢积不满足结合律。

\begin{advexamplebox}{外代数总维数}\label{advA3:ex:exterior-dimension}
证明 $n$ 维空间的外代数满足
\[
 \dim\Lambda(\mathbb R^n)=2^n.
\]

\AQExterior
\end{advexamplebox}
\begin{advsolution}
\AAExterior

$l$ 阶基由 $n$ 个基矢中任选 $l$ 个并按递增次序楔积得到，故
$\dim\Lambda^l=\binom nl$。于是
\[
 \dim\Lambda=\sum_{l=0}^n\binom nl=(1+1)^n=2^n.
\]
这也给出三维的 $1+3+3+1=8$。
\end{advsolution}

\begin{advexamplebox}{霍奇对偶恢复点积}\label{advA3:ex:hodge-dot}
证明对三维矢量 $\vect u,\vect v$，
\[
 \boxed{\star(\vect u\wedge\star\vect v)=\vect u\cdot\vect v}.
\]

\AQHodge
\end{advexamplebox}
\begin{advsolution}
\AAHodge

写 $\star\vect v=\tfrac12v_i\varepsilon_{ijk}\vect e_j\wedge\vect e_k$，则
\[
 \star(\vect u\wedge\star\vect v)
 =\frac12u_lv_i\varepsilon_{ijk}\varepsilon_{ljk}
 =u_lv_i\delta_{li}=u_iv_i.
\]
第二步用了 $\varepsilon_{ijk}\varepsilon_{ljk}=2\delta_{il}$。结果是标量，阶数也由
$1+(3-1)=3$ 再对偶到0阶，检查一致。
\end{advsolution}

\subsection{三维矢量微分}

在每点赋一个标量、矢量或张量分别得到标量场 $f(x)$、矢量场
$\vect F=F_i(x)\vect e_i$ 和张量场。直角坐标固定基中
\[
 \nabla f=(\partial_i f)\vect e_i,
 \quad
 \nabla\vect F=(\partial_iF_j)\vect e_i\vect e_j,
 \quad
 \nabla\cdot\vect F=\partial_iF_i,
 \quad
 \nabla\times\vect F=\varepsilon_{ijk}\partial_iF_j\vect e_k.
\]

从全微分出发，
\[
 \dd f=(\partial_if)\dd x^i=\nabla f\cdot\dd\vect x,
\]
所以梯度是“与任意位移点乘后给出一阶变化”的唯一矢量。继续微分得到 Hessian：
\[
 \dd^2f=(\partial_i\partial_jf)\dd x^i\dd x^j
 =(\dd\vect x\,\dd\vect x):\nabla\nabla f.
\]
矢量场的微分为
\[
 \dd\vect F=(\partial_iF_j)\dd x^i\vect e_j
 =\dd\vect x\cdot\nabla\vect F.
\]
这三式解释了为什么 $\nabla f$ 是矢量、$\nabla\vect F$ 是二阶张量，也解释了张量指标的次序：第一个指标来自求导方向，第二个指标来自场的分量方向。

\begin{advtheorem}{梯度的乘法与链式规则}\label{advA3:thm:gradient-rules}
\[
 \nabla(\lambda f+\mu g)=\lambda\nabla f+\mu\nabla g,
 \quad \nabla(fg)=f\nabla g+g\nabla f,
\]
\[
 \nabla F(g)=F'(g)\nabla g,
 \quad
 \nabla(f\vect F)=f\nabla\vect F+(\nabla f)\vect F,
\]
\[
 \nabla(\vect F\cdot\vect G)
 =(\nabla\vect F)\cdot\vect G+(\nabla\vect G)\cdot\vect F.
\]
\end{advtheorem}
证明逐分量应用普通乘法和链式法则即可。最后两式要保留张量积的指标顺序，不能把 $(\nabla f)\vect F$ 随意换成 $\vect F\nabla f$。

例如最后两式的分量推导为
\[
 \partial_i(fF_j)=f\partial_iF_j+(\partial_if)F_j,
\]
\[
 \partial_i(F_jG_j)=(\partial_iF_j)G_j+(\partial_iG_j)F_j.
\]
第一式右端第二项确实是 $(\nabla f)\vect F$，而非 $\vect F(\nabla f)$；后者的分量是 $F_i\partial_jf$，除非张量恰好对称，否则两者不同。许多矢量恒等式错误都来自把张量积误当成普通数乘，因此这里必须保留指标顺序。

\begin{advexamplebox}{径向函数的梯度与Hessian}\label{advA3:ex:radial-gradients}
令 $r=(x_ix_i)^{1/2}$、$\vect r=x_i\vect e_i$。计算
$\nabla x_i,\nabla r^2,\nabla r,\nabla(1/r),\nabla(\vect r/r^3)$ 以及 $\nabla\nabla r^n$（$r\ne0$）。

\AQRadialHessian
\end{advexamplebox}
\begin{advsolution}
\AARadialHessian

\[
 \boxed{\nabla x_i=\vect e_i},\qquad
 \boxed{\nabla r^2=2\vect r},\qquad
 \boxed{\nabla r=\frac{\vect r}{r}},
\]
\[
 \boxed{\nabla\frac1r=-\frac{\vect r}{r^3}},
 \qquad
 \boxed{\nabla\cdot\frac{\vect r}{r^3}=0\quad(r\ne0)}.
\]
最后一式若把“$\nabla$ 作用于矢量”理解为完整梯度，则
\[
 \nabla\left(\frac{\vect r}{r^3}\right)
 =\frac{\mathsf1}{r^3}-3\frac{\vect r\vect r}{r^5}.
\]
一般地
\[
 \boxed{\nabla\nabla r^n
 =nr^{n-2}\mathsf1+n(n-2)r^{n-4}\vect r\vect r}.
\]
其迹给出 $\nabla^2r^n=n(n+1)r^{n-2}$。$r=0$ 需按分布处理，例如
$\nabla\cdot(\vect r/r^3)=4\pi\delta^{(3)}(\vect r)$。
\end{advsolution}

外导数可写成 $\nabla\wedge(\cdot)=\vect e_i\wedge\partial_i(\cdot)$。
对标量，它就是梯度；对矢量取霍奇对偶得到旋度；对矢量的霍奇对偶再外导并对偶，得到散度。

\begin{advtheorem}{庞加莱引理}\label{advA3:thm:poincare}
外导数恒满足 $\dd^2=0$。在星形区域上，每个次数 $k\ge1$ 的闭微分形式都是恰当形式；对一般区域，这一结论局部成立。零次闭形式只是局部常数，不属于上述正次数结论。特别地
\[
 \nabla\times\nabla f=0,
 \qquad
 \nabla\cdot(\nabla\times\vect F)=0.
\]
\end{advtheorem}
$\dd^2=0$ 直接来自偏导可交换与楔积反对称：
$\partial_i\partial_j=\partial_j\partial_i$，而
$\dd x^i\wedge\dd x^j=-\dd x^j\wedge\dd x^i$。闭形式局部恰当的构造可沿从星形中心到 $x$ 的径向线积分给出。

\begin{advtheorem}{广义斯托克斯定理}\label{advA3:thm:general-stokes}
对有向光滑流形 $M$ 及紧支撑微分形式 $\omega$，
\[
 \boxed{\int_M\dd\omega=\int_{\partial M}\omega}.
\]
\end{advtheorem}
它在坐标小块上由一维微积分基本定理反复使用得到，再用单位分解拼接；内部相邻小块的公共边界因取向相反而抵消，只留下 $\partial M$。

\begin{advtheorem}{散度定理}\label{advA3:thm:divergence-theorem}
\[
 \boxed{\oint_{\partial V}\vect F\cdot\dd\vect S
 =\int_V\nabla\cdot\vect F\,\dd V}.
\]
\end{advtheorem}

\begin{advtheorem}{开尔文--斯托克斯定理}\label{advA3:thm:curl-stokes}
\[
 \boxed{\oint_{\partial S}\vect F\cdot\dd\vect l
 =\int_S(\nabla\times\vect F)\cdot\dd\vect S}.
\]
\end{advtheorem}
两式分别把广义斯托克斯定理应用于 $\star\vect F$ 与矢量场对应的1-形式；边界方向与面法向必须满足右手定向。

\begin{advtheorem}{散度运算公式}\label{advA3:thm:divergence-rules}
\[
 \nabla\cdot(f\vect F)=\nabla f\cdot\vect F+f\nabla\cdot\vect F,
\]
若并矢 $(\vect F\vect G)_{ij}=F_iG_j$ 且散度缩并第一个指标，则
\[
 \nabla\cdot(\vect F\vect G)
 =(\nabla\cdot\vect F)\vect G+(\vect F\cdot\nabla)\vect G.
\]
\[
 \nabla\cdot(\vect F\times\vect G)
 =\vect G\cdot\nabla\times\vect F-\vect F\cdot\nabla\times\vect G,
\]
并且对二阶张量 $\mathsf T$，$(\nabla\cdot\mathsf T)_j=\partial_iT_{ij}$。线性规则逐项成立。
\end{advtheorem}

这些公式都可直接审查指标。例如
\[
 [\nabla\cdot(\vect F\vect G)]_j
 =\partial_i(F_iG_j)=(\partial_iF_i)G_j+F_i\partial_iG_j,
\]
而
\[
 \nabla\cdot(\vect F\times\vect G)
 =\partial_i(\varepsilon_{jki}F_jG_k)
 =G_k\varepsilon_{kij}\partial_iF_j
 -F_j\varepsilon_{jik}\partial_iG_k.
\]
这正是 $\vect G\cdot\nabla\times\vect F-\vect F\cdot\nabla\times\vect G$。指标推导还明确了散度对并矢缩并的是哪个位置，避免把 $(\vect F\cdot\nabla)\vect G$ 错写成 $(\nabla\vect G)\cdot\vect F$ 后再随意转置。

\begin{advtheorem}{旋度与拉普拉斯公式}\label{advA3:thm:curl-rules}
\[
 \nabla\times(f\vect F)=\nabla f\times\vect F+f\nabla\times\vect F,
\]
\[
 \boxed{\nabla\times(\nabla\times\vect F)
 =\nabla(\nabla\cdot\vect F)-\nabla^2\vect F},
\]
\[
 \nabla(\vect F\cdot\vect G)
 =(\vect F\cdot\nabla)\vect G+(\vect G\cdot\nabla)\vect F
 +\vect F\times(\nabla\times\vect G)
 +\vect G\times(\nabla\times\vect F),
\]
\[
 \nabla\times(\vect F\times\vect G)
 =(\vect G\cdot\nabla)\vect F+(\nabla\cdot\vect G)\vect F
 -(\vect F\cdot\nabla)\vect G-(\nabla\cdot\vect F)\vect G.
\]
\end{advtheorem}
证明都可把叉积写成 $\varepsilon_{ijk}$ 并使用
$\varepsilon_{ijk}\varepsilon_{imn}=\delta_{jm}\delta_{kn}-\delta_{jn}\delta_{km}$；这同时避免凭图形直觉误换指标顺序。

以最关键的双旋度为例，第 $i$ 分量为
\[
 [\nabla\times(\nabla\times\vect F)]_i
 =\varepsilon_{jki}\partial_j
 (\varepsilon_{lmk}\partial_lF_m).
\]
缩并两个 $\varepsilon$ 后得到
\[
 (\delta_{il}\delta_{jm}-\delta_{im}\delta_{jl})
 \partial_j\partial_lF_m
 =\partial_i(\partial_mF_m)-\partial_j\partial_jF_i.
\]
因此
$\nabla\times\nabla\times\vect F=\nabla(\nabla\cdot\vect F)-\nabla^2\vect F$。
其余两个含叉积的公式以同一恒等式展开；每项都只使用乘法法则与偏导交换，故在直角坐标固定基中成立。

\begin{advtheorem}{标量、矢量与张量的边界积分形式}\label{advA3:thm:boundary-integrals}
在分段光滑区域上，
\[
 \oint_{\partial V}f\,\dd\vect S=\int_V\nabla f\,\dd V,
 \qquad
 \oint_{\partial V}\dd\vect S\times\vect F
 =\int_V\nabla\times\vect F\,\dd V,
\]
\[
 \oint_{\partial V}\dd\vect S\,\vect F
 =\int_V\nabla\vect F\,\dd V,
 \qquad
 \oint_{\partial V}\dd\vect S\cdot\vect F
 =\int_V\nabla\cdot\vect F\,\dd V,
\]
以及
\[
 \oint_{\partial S}f\,\dd\vect l
 =\int_S\dd\vect S\times\nabla f,
 \qquad
 \oint_{\partial S}\dd\vect l\times\vect F
 =\int_S\bigl[(\nabla\vect F)\cdot\dd\vect S
 -(\nabla\cdot\vect F)\dd\vect S\bigr],
 \]
\[
 \oint_{\partial S}\vect F\cdot\dd\vect l
 =\int_S(\nabla\times\vect F)\cdot\dd\vect S.
\]
\end{advtheorem}
证明：前四式分别把散度定理应用于 $f\vect a$、
$\vect a\cdot(\vect F\times\cdot)$、$T_{ij}$ 的每个分量和 $\vect F$，其中 $\vect a$ 为任意常矢量；任意性保证矢量或张量等式。后两式同理由开尔文--斯托克斯定理逐分量得到。

中间的线积分叉积公式可用任意常矢量 $\vect a$ 核验：
\[
 \vect a\cdot\oint\dd\vect l\times\vect F
 =\oint(\vect F\times\vect a)\cdot\dd\vect l
 =\int_S\nabla\times(\vect F\times\vect a)\cdot\dd\vect S.
\]
由于 $\vect a$ 为常量，
$\nabla\times(\vect F\times\vect a)
=(\vect a\cdot\nabla)\vect F-\vect a(\nabla\cdot\vect F)$，把 $\vect a$ 提出即得原矢量等式。这个证明也固定了两个项的正负号和张量乘法次序。

\begin{advexamplebox}{格林第一、第二恒等式}\label{advA3:ex:green-identities}
证明
\[
 \int_V(\phi\nabla^2\psi+\nabla\phi\cdot\nabla\psi)\dd V
 =\oint_{\partial V}\phi\nabla\psi\cdot\dd\vect S,
\]
\[
 \int_V(\phi\nabla^2\psi-\psi\nabla^2\phi)\dd V
 =\oint_{\partial V}(\phi\nabla\psi-\psi\nabla\phi)\cdot\dd\vect S.
\]

\AQGreen
\end{advexamplebox}
\begin{advsolution}
\AAGreen

乘法法则给出
\[
 \nabla\cdot(\phi\nabla\psi)
 =\nabla\phi\cdot\nabla\psi+\phi\nabla^2\psi.
\]
在 $V$ 上积分并用散度定理即得第一式。再交换 $\phi,\psi$ 得到另一份第一式，两式相减，交叉梯度项相消，得到第二式。若边界是分片光滑的，按每片外法向积分即可；内部人为分界面仍因法向相反而抵消。
\end{advsolution}
